\section{Introduction}

量子色动力学（QCD）是把强子束缚态与其部分子组分——夸克和胶子——联系起来的强相互作用理论。尽管实验上无法直接探测夸克与胶子，
但强子与部分子之间的联系可以通过部分子
分布函数（PDFs）来刻画。

PDFs 描绘了强子结构中与组分夸克、胶子的动量、角动量和自旋相关的方面，在我们理解高能强子散射过程中发挥着核心作用
（例如参见文献
\cite{Brock:1993sz,Sterman:1994ce,Collins:2011zzd}）。通过因子化，
深度非弹性散射和 Drell--Yan 产生等简单散射过程的散射
振幅可以表示为微扰系数与 PDFs 的卷积，后者编码了强子能标下 QCD 的非微扰动力学。

原则上，从 QCD 直接计算 PDFs 将为强子结构提供新的洞见、给出对 QCD 更严格的检验，并降低高能散射实验中的系统不确定度。然而，目前唯一系统的从头计算非微扰 QCD 的方法是格点 QCD，其中 QCD 被构建在离散的欧几里得超立方体上。PDFs 由光前波函数的矩阵元定义，无法直接从欧几里得格点 QCD 确定。目前 PDFs 是通过对大量散射数据的全局分析来确定的（近期若干分析的例子可参见文献
\cite{Arbabifar:2013tma,Ball:2013lla,Jimenez-Delgado:2013boa,
Jimenez-Delgado:2014xza,Dulat:2015mca,
Hou:2016sho,Shahri:2016uzl}）。

因此，格点 QCD 计算转而集中于 PDFs 的前几个 Mellin 矩，它们可以与局域扭度2算符的矩阵元联系起来；这里扭度（twist）是算符的量纲减去自旋。格点调节子破坏了转动对称性，从而诱导出连续时空中本不会混合的格点算符之间的混合。不同质量量纲的扭度2算符之间的混合在格点上引入幂发散混合，使得人们无法提取三个以上的 PDFs 矩
\cite{Detmold:2001dv,Detmold:2003rq}。

最近，人们提出了一种在格点上确定 PDFs 的新途径：借助通常称为准 PDFs（quasi PDFs）的欧几里得对应物
\cite{Ji:2013dva,Ji:2014gla,Lin:2014zya,Gamberg:2014zwa,Ji:2015jwa,
Alexandrou:2015rja,Chen:2016utp}。文献 \cite{Ma:2014jla} 提出了一个类似的框架。
准 PDFs 是在有限核子动量下确定的欧几里得矩阵元。在大欧几里得动量下，准 PDFs 可以通过一种有效理论展开——大动量有效理论（LaMET）——与真实 PDFs 联系起来。

初步的格点计算结果令人鼓舞
\cite{Lin:2014zya,Alexandrou:2015rja}，尽管这两项计算都只采用了单一格距，对其系统不确定度的完整认识仍远未完成。特别是，就其现有形式而言，该途径面临三大挑战：
当前计算资源下核子动量的限制；
对延展欧几里得算符重正化的完整理解；
以及光前 PDFs 与欧几里得准
PDFs 之间的精确关系。

这些困难大体可分为以实际为主或以理论为主两类。第一个挑战，即当前格子上核子动量取值偏低所对应的系统不确定度问题，基本上是一个实际问题。基于旁观者双夸克模型的研究
\cite{Gamberg:2014zwa} 表明，计算资源的适度改进以及针对大动量核子定制的新算法 \cite{Bali:2016lva}，很可能解决这一困难，其精度至少足以对实验数据难以覆盖的参数区域的全局 PDF 分析作出贡献。下文我们不再考虑这些实际
困难，转而专注于准 PDFs 的理论
方面。

我们通过提出一种在格点上计算准 PDFs 的新方法来应对其中一个理论挑战：在该方法中，格点自由度经由梯度流被涂抹
\cite{Narayanan:2006rf,Luscher:2011bx,Luscher:2013cpa}。使用带环（ringed）费米子——它们不需要任何乘法性波函数重正化
\cite{Makino:2014taa,Hieda:2016lly}——相应的格点矩阵元在连续极限下保持有限。该方法回避了定义准 PDF 的 Wilson 线算符所伴有的幂发散问题。准 PDFs 的重正化曾借助重夸克有效理论的视角加以审视 \cite{Ji:2015jwa}；近来亦有人提出了消除这一幂发散的反项步骤
\cite{Chen:2016fxx,Ishikawa:2016znu}，但这两种方法都尚未在微扰论二圈以外得到确立。


本文考察涂抹准
PDF 与光前 PDF 之间的关系，聚焦于流时间与核子动量设定的长度尺度相比很小的极限，此时核子动量足够大以致高扭度效应可以忽略。在这一极限下，我们通过小流时间展开把涂抹准 PDF 的矩用光前 PDF 的矩表出。
我们方法的首要优点是：用格点 QCD 确定的非微扰矩阵元具有有限的连续极限。由流时间调节的涂抹准 PDF 与光前 PDF（比如取 $\ms$
方案）之间的匹配可在连续时空中进行，并且与非微扰格点计算的细节无关。

我们先在第~\ref{sec:defs}~节重温光前 PDFs 与准 PDFs 的定义。随后在第~\ref{sec:spdf}~节分析光前 PDFs 与准 PDFs 之间的关系，并在第~\ref{sec:match}~节推导匹配核的演化方程。最后我们在第~\ref{sec:concs}~节给出总结与结论。


